% !TEX root = ../main.tex
\chapter{Related Work}
\label{ch:relatedwork}

% =====================================================================
% 骨架來源：docs/thesis/thesis-outline.md §2（background）+ §11（related work）
% 兩者合併成樣板原本的第 2 章。前四節是背景與問題刻畫，後七節是文獻。
%
% 相依：docs/thesis/related-work-v1-thought.md（已離線審完）
%
% 【使用該產出前必讀】它的「已驗證」實際只等於「連結打得開」：
%   - 至少 12 篇有同儕審查場所的論文被標成 arXiv preprint（場所欄要用 DBLP 重掃）
%   - 六篇 2026 的論文無法離線核對，其中 arXiv 2607.08010（Amazon tool-making）
%     一篇獨力推翻桶 1 的 gap，核對完成前不得據此改任何內容
%   - 桶 3（LLM + PBE）時間軸停在 2023，2024–2026 全空白，這是唯一非補不可的洞
% =====================================================================

\section{Template-Shaped Tasks}
\label{sec:rel:tasks}

% TODO 任務定義：給定模板 T 與一組值 V，產生實例 D = f(T, V)。
%      難點在 f 的定位（locator）部分：哪一個儲存格／bookmark／文字對應哪一個語意欄位。

\section{Document Formats and Representation Mismatch}
\label{sec:rel:formats}

% TODO docx / xlsx 的資料模型：bookmark、content control、具名範圍、合併儲存格、
%      公式 vs 快取值。以及 DWG 的 entity graph、layer、block/attribute、xdata。
% TODO 表徵不匹配：語意存在於佈局關係，序列化成 token 會摧毀這些結構。
%      【降階處理】這是動機的一部分但不能當主論點，會被一句「下一代模型就解決了」擊破。
%      正確的推論是：用工具提供結構化 dump，並且只在編譯期讀一次。

\section{Failure Modes}
\label{sec:rel:failures}

% TODO 本節是後面 oracle 與 gate 設計的依據，第 3/4 章要有一張對照表回指這裡：
%      - 沉默失敗（值錯、位置對）
%      - 定位漂移（模板改版後 locator 指到別的地方）
%      - 覆蓋不足（欄位根本沒被規格涵蓋 -> 保持模板預設值，最難察覺）
%      - 格式／單位錯誤
%      - negative obligation 違反（該空的沒清空）—— 新增，來自 benchmark S5
%      - 範圍外延錯誤（聚合多算或少算一列，數字有理有據但錯）—— 新增，來自 ADR 0002

\section{Governance Requirements}
\label{sec:rel:governance}

% TODO 可重現、可稽核、可簽核、可回溯。這是純 agent 方案的結構性障礙。
% 【措辭修正】不要寫「法規普遍要求 deterministic output」——找不到學術來源。
%      改寫成 traceability / reproducibility / change control。

\section{LLMs for Document and Engineering-Drawing Understanding}
\label{sec:rel:understanding}

\section{Document Automation, RPA, and Template Filling}
\label{sec:rel:rpa}

% TODO 三分法在這裡展開，但結論句已經在 Introduction 出現過：
%      RPA = deterministic but manual / agents = automatic but stochastic /
%      ours = automatic derivation + deterministic execution

\section{Programming by Example and Program Synthesis}
\label{sec:rel:pbe}

% TODO 【這一節是第二危險的鄰居，也是產出裡唯一非補不可的洞：2024–2026 全空白】
%      FlashFill++ 一類的工作是最直接的比較對象。
%      區別點：範圍聚合（S3）與 negative obligation（S5）是兩個無法化約成單欄 PBE 的案例。

\section{Tool-Making, Skill Libraries, and Amortized LLM Computation}
\label{sec:rel:toolmaking}

% TODO 【第一名危險鄰居】LATM、Voyager skill library、TroVE、CRAFT、Agent Skills，
%      以及待核對的 arXiv 2607.08010（Amazon tool-making）。
%      Gap：這些工作把產生的工具存下來就結束——沒有可執行 oracle、沒有覆蓋率概念、
%      沒有回歸與版本治理。而在有責任的交付物上，這三者才決定能否上線。

\section{Execution-Feedback Program Repair}
\label{sec:rel:repair}

% TODO Self-Debug、Reflexion、AlphaCodium。
%      Gap：它們的 oracle 多來自現成單元測試；本工作的 oracle 必須從文件比對建構，
%      而且要處理「沉默失敗」而非「程式當掉」。

\section{Verification and Guardrails for LLM Output}
\label{sec:rel:verification}

% TODO LLM-as-judge 的侷限、property-based / metamorphic testing。
%      支撐第 4 章的「要 oracle 不要 judge」與第 5 章評測器不得問 LLM 的決定。

\section{Cost-Aware Agent Evaluation}
\label{sec:rel:evaluation}

% TODO Gap：多數 agent 評測只報準確率，不報變異度與攤提成本，
%      而這兩項正是生產環境的決策依據。這一節直接支撐第 5 章的指標選擇。
